\documentclass[11pt,a4paper]{article}

\usepackage[utf8]{inputenc}
\usepackage[T1]{fontenc}
\usepackage[french]{babel}
\usepackage[a4paper,margin=2.3cm]{geometry}
\usepackage{amsmath,amssymb,amsthm,mathtools}
\usepackage{lmodern}
\usepackage{enumitem}
\usepackage{hyperref}
\usepackage{xcolor}
\usepackage{fancyhdr}
\usepackage{stmaryrd}
\usepackage{mathrsfs}
\usepackage{array}

\hypersetup{
    colorlinks=true,
    linkcolor=black,
    urlcolor=blue,
    citecolor=black
}

\setlength{\parindent}{0pt}
\setlength{\parskip}{0.6em}

\pagestyle{fancy}
\fancyhf{}
\fancyhead[L]{Algèbre linéaire --- Chapitre 4}
\fancyhead[R]{Calcul matriciel}
\fancyfoot[C]{\thepage}

% Environnements
\newtheorem{definition}{Définition}[section]
\newtheorem{theoreme}[definition]{Théorème}
\newtheorem{proposition}[definition]{Proposition}
\newtheorem{corollaire}[definition]{Corollaire}
\newtheorem{lemme}[definition]{Lemme}
\theoremstyle{remark}
\newtheorem{remarque}[definition]{Remarque}
\newtheorem{exemple}[definition]{Exemple}
\newtheorem*{methode}{Méthode}
\newtheorem*{idee}{Idée}

% Commandes
\newcommand{\R}{\mathbb{R}}
\newcommand{\C}{\mathbb{C}}
\newcommand{\K}{\mathbb{K}}
\newcommand{\Mnp}{M_{n,p}(\K)}
\newcommand{\Mnn}{M_n(\K)}
\newcommand{\GLn}{GL_n(\K)}
\newcommand{\rg}{\mathrm{rg}}
\newcommand{\tr}{\mathrm{tr}}

\begin{document}

\begin{center}
    {\LARGE \textbf{Chapitre 4 --- Calcul matriciel}}\\[0.5em]
    {\large Manipuler efficacement les écritures linéaires}
\end{center}

\hrule
\vspace{1em}

\tableofcontents

\newpage

\section{Pourquoi les matrices sont-elles fondamentales ?}

Les systèmes linéaires étudiés au chapitre précédent font apparaître naturellement des tableaux de coefficients. Ces tableaux sont les \textbf{matrices}. Elles permettent :
\begin{itemize}
    \item de condenser l’écriture d’un système linéaire ;
    \item d’organiser les calculs ;
    \item de représenter des transformations linéaires ;
    \item de passer d’un calcul sur des vecteurs à un calcul plus structuré.
\end{itemize}

\begin{idee}
Une matrice est avant tout un \textbf{outil d’organisation} de l’information linéaire.
\end{idee}

\begin{remarque}
Le calcul matriciel n’est pas seulement un langage pratique : il possède ses propres règles, qu’il faut maîtriser soigneusement. En particulier, le produit matriciel n’a pas les mêmes propriétés que le produit de nombres réels.
\end{remarque}

\section{Définition d’une matrice}

\subsection{Définition}

\begin{definition}
Soient \(n,p\in \N^\ast\).  
Une \textbf{matrice à \(n\) lignes et \(p\) colonnes} à coefficients dans \(\K\) est un tableau rectangulaire de la forme
\[
A=(a_{ij})_{1\le i\le n,\ 1\le j\le p},
\]
où chaque coefficient \(a_{ij}\in \K\).
\end{definition}

On note
\[
M_{n,p}(\K)
\]
l’ensemble des matrices à \(n\) lignes et \(p\) colonnes à coefficients dans \(\K\).

\begin{remarque}
Le coefficient \(a_{ij}\) est situé à l’intersection de la \(i\)-ème ligne et de la \(j\)-ème colonne.
\end{remarque}

\begin{exemple}
\[
A=
\begin{pmatrix}
1 & 2 & -1\\
0 & 3 & 4
\end{pmatrix}
\in M_{2,3}(\R).
\]
Ici :
\[
a_{11}=1,\quad a_{12}=2,\quad a_{13}=-1,\quad a_{21}=0,\quad a_{22}=3,\quad a_{23}=4.
\]
\end{exemple}

\subsection{Taille d’une matrice}

\begin{definition}
Si \(A\in M_{n,p}(\K)\), on dit que \(A\) est de taille \(n\times p\).
\end{definition}

\begin{exemple}
\[
\begin{pmatrix}
1 & 0\\
2 & 3\\
-1 & 4
\end{pmatrix}
\]
est une matrice de taille \(3\times 2\).
\end{exemple}

\subsection{Matrices carrées}

\begin{definition}
Une matrice est dite \textbf{carrée} lorsque son nombre de lignes est égal à son nombre de colonnes, c’est-à-dire lorsqu’elle appartient à \(M_n(\K)=M_{n,n}(\K)\).
\end{definition}

\begin{exemple}
\[
\begin{pmatrix}
1 & 2\\
3 & 4
\end{pmatrix}
\in M_2(\R)
\]
est une matrice carrée d’ordre \(2\).
\end{exemple}

\section{Égalité de matrices}

\begin{definition}
Deux matrices \(A=(a_{ij})\) et \(B=(b_{ij})\) sont égales si et seulement si :
\begin{itemize}
    \item elles ont la même taille ;
    \item tous leurs coefficients correspondants sont égaux :
    \[
    \forall i,\forall j,\ a_{ij}=b_{ij}.
    \]
\end{itemize}
\end{definition}

\begin{exemple}
Si
\[
\begin{pmatrix}
x & 1\\
2 & y
\end{pmatrix}
=
\begin{pmatrix}
3 & 1\\
2 & -4
\end{pmatrix},
\]
alors
\[
x=3\qquad \text{et} \qquad y=-4.
\]
\end{exemple}

\section{Addition et multiplication par un scalaire}

\subsection{Addition}

\begin{definition}
Soient \(A=(a_{ij})\) et \(B=(b_{ij})\) deux matrices de même taille \(n\times p\).  
On définit leur somme par
\[
A+B=(a_{ij}+b_{ij})_{1\le i\le n,\ 1\le j\le p}.
\]
\end{definition}

\begin{exemple}
\[
\begin{pmatrix}
1 & 2\\
3 & 4
\end{pmatrix}
+
\begin{pmatrix}
0 & -1\\
5 & 2
\end{pmatrix}
=
\begin{pmatrix}
1 & 1\\
8 & 6
\end{pmatrix}.
\]
\end{exemple}

\subsection{Multiplication par un scalaire}

\begin{definition}
Soit \(\lambda\in \K\) et \(A=(a_{ij})\in M_{n,p}(\K)\).  
On définit
\[
\lambda A=(\lambda a_{ij})_{1\le i\le n,\ 1\le j\le p}.
\]
\end{definition}

\begin{exemple}
\[
2\begin{pmatrix}
1 & -3\\
0 & 4
\end{pmatrix}
=
\begin{pmatrix}
2 & -6\\
0 & 8
\end{pmatrix}.
\]
\end{exemple}

\subsection{Structure d’espace vectoriel}

\begin{proposition}
L’ensemble \(M_{n,p}(\K)\), muni de l’addition et de la multiplication par un scalaire définies coefficient par coefficient, est un espace vectoriel sur \(\K\).
\end{proposition}

\begin{proof}
Les propriétés d’espace vectoriel découlent directement des propriétés des opérations dans \(\K\), appliquées coefficient par coefficient.
\end{proof}

\begin{remarque}
On retrouve ici un exemple fondamental d’espace vectoriel de dimension finie.
\end{remarque}

\section{Matrices particulières}

\subsection{Matrice nulle}

\begin{definition}
La \textbf{matrice nulle} de \(M_{n,p}(\K)\), notée \(0_{n,p}\) ou simplement \(0\) lorsqu’il n’y a pas d’ambiguïté, est la matrice dont tous les coefficients sont nuls.
\end{definition}

\begin{exemple}
\[
0_{2,3}=
\begin{pmatrix}
0 & 0 & 0\\
0 & 0 & 0
\end{pmatrix}.
\]
\end{exemple}

\subsection{Matrice identité}

\begin{definition}
Dans \(M_n(\K)\), la \textbf{matrice identité} d’ordre \(n\), notée \(I_n\), est la matrice définie par
\[
(I_n)_{ij}=
\begin{cases}
1 & \text{si } i=j,\\
0 & \text{si } i\neq j.
\end{cases}
\]
\end{definition}

\begin{exemple}
\[
I_2=
\begin{pmatrix}
1 & 0\\
0 & 1
\end{pmatrix},
\qquad
I_3=
\begin{pmatrix}
1 & 0 & 0\\
0 & 1 & 0\\
0 & 0 & 1
\end{pmatrix}.
\]
\end{exemple}

\subsection{Matrices lignes et colonnes}

\begin{definition}
Une matrice à une seule ligne est appelée \textbf{matrice ligne}.  
Une matrice à une seule colonne est appelée \textbf{matrice colonne}.
\end{definition}

\begin{exemple}
\[
\begin{pmatrix}
1 & -2 & 4
\end{pmatrix}
\]
est une matrice ligne, tandis que
\[
\begin{pmatrix}
1\\
-2\\
4
\end{pmatrix}
\]
est une matrice colonne.
\end{exemple}

\subsection{Matrices diagonales, scalaires, triangulaires}

\begin{definition}
Une matrice carrée \(A=(a_{ij})\in M_n(\K)\) est :
\begin{itemize}
    \item \textbf{diagonale} si \(a_{ij}=0\) pour \(i\neq j\) ;
    \item \textbf{scalaire} si elle est diagonale et si tous les coefficients diagonaux sont égaux ;
    \item \textbf{triangulaire supérieure} si \(a_{ij}=0\) pour \(i>j\) ;
    \item \textbf{triangulaire inférieure} si \(a_{ij}=0\) pour \(i<j\).
\end{itemize}
\end{definition}

\begin{exemple}
\[
\begin{pmatrix}
1 & 0 & 0\\
0 & -2 & 0\\
0 & 0 & 5
\end{pmatrix}
\]
est diagonale.

\[
3I_3=
\begin{pmatrix}
3 & 0 & 0\\
0 & 3 & 0\\
0 & 0 & 3
\end{pmatrix}
\]
est scalaire.

\[
\begin{pmatrix}
1 & 2 & 3\\
0 & 4 & 5\\
0 & 0 & 6
\end{pmatrix}
\]
est triangulaire supérieure.
\end{exemple}

\section{Produit matriciel}

\subsection{Définition}

\begin{definition}
Soient
\[
A=(a_{ij})\in M_{n,p}(\K)
\qquad \text{et} \qquad
B=(b_{jk})\in M_{p,q}(\K).
\]
On définit leur produit \(AB\in M_{n,q}(\K)\) par
\[
(AB)_{ik}=\sum_{j=1}^p a_{ij}b_{jk}.
\]
\end{definition}

\begin{remarque}
Le produit \(AB\) n’est défini que si le nombre de colonnes de \(A\) est égal au nombre de lignes de \(B\).
\end{remarque}

\begin{idee}
Le coefficient \((i,k)\) du produit s’obtient en faisant le produit scalaire de la \(i\)-ème ligne de \(A\) avec la \(k\)-ème colonne de \(B\).
\end{idee}

\subsection{Exemple de calcul}

\begin{exemple}
Soient
\[
A=
\begin{pmatrix}
1 & 2\\
3 & 4
\end{pmatrix},
\qquad
B=
\begin{pmatrix}
5 & 6\\
7 & 8
\end{pmatrix}.
\]
Alors
\[
AB=
\begin{pmatrix}
1\cdot 5+2\cdot 7 & 1\cdot 6+2\cdot 8\\
3\cdot 5+4\cdot 7 & 3\cdot 6+4\cdot 8
\end{pmatrix}
=
\begin{pmatrix}
19 & 22\\
43 & 50
\end{pmatrix}.
\]
\end{exemple}

\begin{exemple}
Soient
\[
A=
\begin{pmatrix}
1 & 2 & 3\\
0 & -1 & 4
\end{pmatrix}
\in M_{2,3}(\R),
\qquad
B=
\begin{pmatrix}
1\\
2\\
-1
\end{pmatrix}
\in M_{3,1}(\R).
\]
Alors
\[
AB=
\begin{pmatrix}
1+4-3\\
0-2-4
\end{pmatrix}
=
\begin{pmatrix}
2\\
-6
\end{pmatrix}.
\]
\end{exemple}

\subsection{Propriétés du produit matriciel}

\begin{proposition}
Soient \(A,B,C\) des matrices de tailles compatibles, et \(\lambda\in\K\). Alors :
\begin{enumerate}[label=\arabic*)]
    \item \(A(BC)=(AB)C\) ;
    \item \(A(B+C)=AB+AC\) ;
    \item \((A+B)C=AC+BC\) ;
    \item \(\lambda(AB)=(\lambda A)B=A(\lambda B)\) ;
    \item \(I_nA=A\) et \(AI_p=A\) si \(A\in M_{n,p}(\K)\).
\end{enumerate}
\end{proposition}

\begin{proof}
Les preuves reposent sur le calcul des coefficients des matrices.

Par exemple, pour l’associativité, si les tailles sont compatibles :
\[
((AB)C)_{ik}
=
\sum_{j} (AB)_{ij}c_{jk}
=
\sum_j \left(\sum_\ell a_{i\ell}b_{\ell j}\right)c_{jk}.
\]
En réorganisant les sommes,
\[
((AB)C)_{ik}
=
\sum_\ell a_{i\ell}\left(\sum_j b_{\ell j}c_{jk}\right)
=
\sum_\ell a_{i\ell}(BC)_{\ell k}
=
(A(BC))_{ik}.
\]
Les autres propriétés se démontrent de façon analogue.
\end{proof}

\begin{remarque}
L’associativité permet d’écrire \(ABC\) sans parenthèses dès que le produit a un sens.
\end{remarque}

\subsection{Non-commutativité}

\begin{proposition}
En général,
\[
AB\neq BA.
\]
\end{proposition}

\begin{exemple}
Posons
\[
A=
\begin{pmatrix}
1 & 1\\
0 & 0
\end{pmatrix},
\qquad
B=
\begin{pmatrix}
1 & 0\\
1 & 0
\end{pmatrix}.
\]
Alors
\[
AB=
\begin{pmatrix}
2 & 0\\
0 & 0
\end{pmatrix},
\qquad
BA=
\begin{pmatrix}
1 & 1\\
1 & 1
\end{pmatrix}.
\]
Donc
\[
AB\neq BA.
\]
\end{exemple}

\begin{remarque}
C’est l’une des différences majeures avec le calcul réel. Il faut toujours faire attention à l’ordre des facteurs.
\end{remarque}

\section{Produit par blocs ligne-colonne}

\begin{proposition}
Si
\[
A=
\begin{pmatrix}
L_1\\
\vdots\\
L_n
\end{pmatrix}
\in M_{n,p}(\K),
\qquad
B=
\begin{pmatrix}
C_1 & \cdots & C_q
\end{pmatrix}
\in M_{p,q}(\K),
\]
où les \(L_i\) sont les lignes de \(A\) et les \(C_j\) les colonnes de \(B\), alors
\[
AB=
\big(L_iC_j\big)_{1\le i\le n,\ 1\le j\le q}.
\]
\end{proposition}

\begin{remarque}
Autrement dit, chaque coefficient de \(AB\) est obtenu en multipliant une ligne de \(A\) par une colonne de \(B\).
\end{remarque}

\section{Puissances d’une matrice carrée}

\begin{definition}
Soit \(A\in M_n(\K)\). On définit, pour tout entier \(k\in \N\),
\[
A^0=I_n,\qquad A^1=A,\qquad A^{k+1}=A^kA.
\]
\end{definition}

\begin{remarque}
Les puissances d’une matrice n’ont de sens que pour les matrices carrées.
\end{remarque}

\begin{proposition}
Pour \(A\in M_n(\K)\) et \(p,q\in\N\),
\[
A^pA^q=A^{p+q}.
\]
\end{proposition}

\begin{proof}
La propriété se démontre par récurrence, en utilisant l’associativité du produit matriciel.
\end{proof}

\begin{exemple}
Si
\[
A=
\begin{pmatrix}
1 & 1\\
0 & 1
\end{pmatrix},
\]
alors
\[
A^2=
\begin{pmatrix}
1 & 2\\
0 & 1
\end{pmatrix},
\qquad
A^3=
\begin{pmatrix}
1 & 3\\
0 & 1
\end{pmatrix}.
\]
\end{exemple}

\section{Transposée}

\subsection{Définition}

\begin{definition}
Soit \(A=(a_{ij})\in M_{n,p}(\K)\).  
La \textbf{transposée} de \(A\), notée \(A^T\), est la matrice de \(M_{p,n}(\K)\) définie par
\[
(A^T)_{ij}=a_{ji}.
\]
\end{definition}

\begin{idee}
Transposer une matrice revient à échanger ses lignes et ses colonnes.
\end{idee}

\begin{exemple}
\[
A=
\begin{pmatrix}
1 & 2 & 3\\
4 & 5 & 6
\end{pmatrix}
\quad \Longrightarrow \quad
A^T=
\begin{pmatrix}
1 & 4\\
2 & 5\\
3 & 6
\end{pmatrix}.
\]
\end{exemple}

\subsection{Propriétés}

\begin{proposition}
Pour toutes matrices \(A,B\) de tailles compatibles et tout scalaire \(\lambda\in\K\),
\begin{enumerate}[label=\arabic*)]
    \item \((A^T)^T=A\) ;
    \item \((A+B)^T=A^T+B^T\) ;
    \item \((\lambda A)^T=\lambda A^T\) ;
    \item \((AB)^T=B^TA^T\).
\end{enumerate}
\end{proposition}

\begin{proof}
Les trois premières propriétés découlent immédiatement de la définition coefficient par coefficient.

Pour la quatrième, si \(A\in M_{n,p}(\K)\) et \(B\in M_{p,q}(\K)\), alors
\[
((AB)^T)_{ij}=(AB)_{ji}=\sum_{k=1}^p a_{jk}b_{ki}.
\]
Or
\[
(B^TA^T)_{ij}=\sum_{k=1}^p (B^T)_{ik}(A^T)_{kj}
=\sum_{k=1}^p b_{ki}a_{jk}.
\]
Comme les produits scalaires commutent dans \(\K\),
\[
((AB)^T)_{ij}=(B^TA^T)_{ij}.
\]
Donc
\[
(AB)^T=B^TA^T.
\]
\end{proof}

\begin{remarque}
Il faut bien retenir l’inversion de l’ordre :
\[
(AB)^T\neq A^TB^T \quad \text{en général.}
\]
\end{remarque}

\subsection{Matrices symétriques}

\begin{definition}
Une matrice carrée \(A\in M_n(\K)\) est dite \textbf{symétrique} si
\[
A^T=A.
\]
\end{definition}

\begin{exemple}
\[
\begin{pmatrix}
1 & 2 & 3\\
2 & 0 & -1\\
3 & -1 & 4
\end{pmatrix}
\]
est symétrique.
\end{exemple}

\section{Matrices inversibles}

\subsection{Définition}

\begin{definition}
Soit \(A\in M_n(\K)\).  
On dit que \(A\) est \textbf{inversible} s’il existe une matrice \(B\in M_n(\K)\) telle que
\[
AB=BA=I_n.
\]
Cette matrice \(B\) est appelée l’\textbf{inverse} de \(A\), et notée \(A^{-1}\).
\end{definition}

\begin{remarque}
Seules les matrices carrées peuvent être inversibles.
\end{remarque}

\subsection{Unicité de l’inverse}

\begin{proposition}
Si une matrice carrée \(A\) est inversible, alors son inverse est unique.
\end{proposition}

\begin{proof}
Supposons que \(B\) et \(C\) soient deux inverses de \(A\). Alors
\[
AB=BA=I_n
\qquad \text{et} \qquad
AC=CA=I_n.
\]
On calcule :
\[
B=BI_n=B(AC)=(BA)C=I_nC=C.
\]
Donc l’inverse est unique.
\end{proof}

\subsection{Exemples}

\begin{exemple}
La matrice identité \(I_n\) est inversible, et
\[
I_n^{-1}=I_n.
\]
\end{exemple}

\begin{exemple}
Si \(\lambda\in\K^\ast\), alors la matrice scalaire \(\lambda I_n\) est inversible et
\[
(\lambda I_n)^{-1}=\frac{1}{\lambda}I_n.
\]
\end{exemple}

\begin{exemple}
Pour
\[
A=
\begin{pmatrix}
a & 0\\
0 & b
\end{pmatrix},
\]
si \(a\neq 0\) et \(b\neq 0\), alors
\[
A^{-1}=
\begin{pmatrix}
\frac1a & 0\\
0 & \frac1b
\end{pmatrix}.
\]
\end{exemple}

\subsection{Produit de matrices inversibles}

\begin{proposition}
Si \(A\) et \(B\) sont deux matrices inversibles de \(M_n(\K)\), alors \(AB\) est inversible et
\[
(AB)^{-1}=B^{-1}A^{-1}.
\]
\end{proposition}

\begin{proof}
On calcule :
\[
(AB)(B^{-1}A^{-1})=A(BB^{-1})A^{-1}=AI_nA^{-1}=I_n,
\]
et
\[
(B^{-1}A^{-1})(AB)=B^{-1}(A^{-1}A)B=B^{-1}I_nB=I_n.
\]
Donc \(B^{-1}A^{-1}\) est bien l’inverse de \(AB\).
\end{proof}

\begin{remarque}
L’ordre est inversé, comme pour la transposée.
\end{remarque}

\subsection{Inverse de la transposée}

\begin{proposition}
Si \(A\) est inversible, alors \(A^T\) est inversible et
\[
(A^T)^{-1}=(A^{-1})^T.
\]
\end{proposition}

\begin{proof}
Comme
\[
AA^{-1}=I_n,
\]
en transposant :
\[
(A^{-1})^TA^T=I_n.
\]
De même, à partir de
\[
A^{-1}A=I_n,
\]
on obtient
\[
A^T(A^{-1})^T=I_n.
\]
Donc \((A^{-1})^T\) est l’inverse de \(A^T\).
\end{proof}

\section{Calcul explicite de l’inverse en petite taille}

\subsection{Cas \(2\times 2\)}

\begin{proposition}
Soit
\[
A=
\begin{pmatrix}
a & b\\
c & d
\end{pmatrix}.
\]
Si
\[
ad-bc\neq 0,
\]
alors \(A\) est inversible et
\[
A^{-1}=\frac{1}{ad-bc}
\begin{pmatrix}
d & -b\\
-c & a
\end{pmatrix}.
\]
\end{proposition}

\begin{proof}
On vérifie directement que
\[
A\cdot
\begin{pmatrix}
d & -b\\
-c & a
\end{pmatrix}
=
\begin{pmatrix}
ad-bc & 0\\
0 & ad-bc
\end{pmatrix}
=(ad-bc)I_2.
\]
Si \(ad-bc\neq 0\), on divise par cette quantité et on obtient bien l’inverse.
\end{proof}

\begin{exemple}
Pour
\[
A=
\begin{pmatrix}
1 & 2\\
3 & 4
\end{pmatrix},
\]
on a
\[
ad-bc=4-6=-2\neq 0,
\]
donc
\[
A^{-1}
=
-\frac12
\begin{pmatrix}
4 & -2\\
-3 & 1
\end{pmatrix}
=
\begin{pmatrix}
-2 & 1\\
\frac32 & -\frac12
\end{pmatrix}.
\]
\end{exemple}

\subsection{Méthode par résolution de systèmes}

\begin{remarque}
Pour calculer l’inverse d’une matrice \(A\), on peut aussi résoudre les systèmes
\[
AX=E_1,\quad AX=E_2,\quad \dots,\quad AX=E_n,
\]
où \(E_1,\dots,E_n\) sont les colonnes de la base canonique. Les solutions forment alors les colonnes de \(A^{-1}\).
\end{remarque}

\section{Opérations élémentaires sur les lignes et matrices élémentaires}

\subsection{Rappel}

\begin{remarque}
Au chapitre précédent, on a utilisé des opérations élémentaires sur les lignes d’un système. Ces opérations peuvent être vues comme des multiplications à gauche par certaines matrices particulières.
\end{remarque}

\subsection{Matrices élémentaires}

\begin{definition}
On appelle \textbf{matrices élémentaires} les matrices obtenues à partir de \(I_n\) en effectuant une unique opération élémentaire sur les lignes :
\begin{itemize}
    \item échange de deux lignes ;
    \item multiplication d’une ligne par un scalaire non nul ;
    \item ajout à une ligne d’un multiple d’une autre.
\end{itemize}
\end{definition}

\begin{proposition}
Multiplier une matrice \(A\) à gauche par une matrice élémentaire revient à effectuer sur \(A\) l’opération élémentaire de lignes correspondante.
\end{proposition}

\begin{remarque}
C’est le fondement matriciel de la méthode du pivot de Gauss.
\end{remarque}

\section{Lien entre matrices et systèmes linéaires}

\subsection{Écriture matricielle}

\begin{definition}
Le système linéaire
\[
\begin{cases}
a_{11}x_1+\cdots+a_{1n}x_n=b_1,\\
\hspace{2.2cm}\vdots\\
a_{p1}x_1+\cdots+a_{pn}x_n=b_p
\end{cases}
\]
s’écrit matriciellement
\[
AX=B,
\]
où
\[
A=(a_{ij})\in M_{p,n}(\K),
\qquad
X=
\begin{pmatrix}
x_1\\
\vdots\\
x_n
\end{pmatrix},
\qquad
B=
\begin{pmatrix}
b_1\\
\vdots\\
b_p
\end{pmatrix}.
\]
\end{definition}

\begin{exemple}
Le système
\[
\begin{cases}
x+2y-z=3,\\
2x-y+4z=1
\end{cases}
\]
s’écrit
\[
\begin{pmatrix}
1 & 2 & -1\\
2 & -1 & 4
\end{pmatrix}
\begin{pmatrix}
x\\
y\\
z
\end{pmatrix}
=
\begin{pmatrix}
3\\
1
\end{pmatrix}.
\]
\end{exemple}

\subsection{Cas d’un système carré}

\begin{proposition}
Si \(A\in M_n(\K)\) est inversible, alors le système
\[
AX=B
\]
admet une unique solution, donnée par
\[
X=A^{-1}B.
\]
\end{proposition}

\begin{proof}
Si \(AX=B\), en multipliant à gauche par \(A^{-1}\), on obtient
\[
X=A^{-1}B.
\]
Cette expression fournit une solution.

Réciproquement, cette solution est unique car toute solution doit satisfaire cette égalité.
\end{proof}

\begin{remarque}
L’inversibilité d’une matrice carrée garantit donc l’existence et l’unicité de la solution pour tout second membre.
\end{remarque}

\section{Matrice d’un vecteur et conventions}

\begin{remarque}
En algèbre linéaire, un vecteur de \(\K^n\) est souvent identifié à une matrice colonne de taille \(n\times 1\). C’est cette convention qui permet d’écrire les systèmes sous la forme
\[
AX=B.
\]
\end{remarque}

\begin{exemple}
Le vecteur \((x,y,z)\in\R^3\) sera souvent écrit
\[
\begin{pmatrix}
x\\
y\\
z
\end{pmatrix}.
\]
\end{exemple}

\section{Premiers résultats utiles sur les matrices carrées}

\subsection{Simplification}

\begin{remarque}
Contrairement aux nombres réels, on ne peut pas toujours simplifier dans un produit matriciel. Par exemple, de
\[
AB=AC
\]
on ne peut pas conclure en général que
\[
B=C.
\]
\end{remarque}

\begin{proposition}
Si \(A\in M_n(\K)\) est inversible et si
\[
AB=AC,
\]
alors
\[
B=C.
\]
De même, si
\[
BA=CA,
\]
alors
\[
B=C.
\]
\end{proposition}

\begin{proof}
Si \(AB=AC\), on multiplie à gauche par \(A^{-1}\) :
\[
A^{-1}AB=A^{-1}AC
\Longrightarrow
B=C.
\]
Le second cas se traite en multipliant à droite par \(A^{-1}\).
\end{proof}

\subsection{Produit nul}

\begin{remarque}
Le produit matriciel n’a pas la propriété de produit nul des réels :
\[
AB=0
\]
n’implique pas en général que
\[
A=0 \quad \text{ou} \quad B=0.
\]
\end{remarque}

\begin{exemple}
Posons
\[
A=
\begin{pmatrix}
1 & 0\\
0 & 0
\end{pmatrix},
\qquad
B=
\begin{pmatrix}
0 & 0\\
1 & 0
\end{pmatrix}.
\]
Alors
\[
AB=
\begin{pmatrix}
0 & 0\\
0 & 0
\end{pmatrix},
\]
mais \(A\neq 0\) et \(B\neq 0\).
\end{exemple}

\section{Dimension de l’espace des matrices}

\begin{proposition}
L’espace vectoriel \(M_{n,p}(\K)\) est de dimension
\[
np.
\]
\end{proposition}

\begin{proof}
Pour \(1\le i\le n\) et \(1\le j\le p\), on note \(E_{ij}\) la matrice dont tous les coefficients sont nuls sauf celui en position \((i,j)\), qui vaut \(1\).

Toute matrice \(A=(a_{ij})\in M_{n,p}(\K)\) s’écrit
\[
A=\sum_{i=1}^n\sum_{j=1}^p a_{ij}E_{ij}.
\]
Donc la famille \((E_{ij})\) engendre \(M_{n,p}(\K)\).

Elle est libre : si
\[
\sum_{i=1}^n\sum_{j=1}^p \lambda_{ij}E_{ij}=0,
\]
alors tous les coefficients de la matrice nulle sont nuls, donc
\[
\lambda_{ij}=0
\]
pour tous \(i,j\).

Cette famille est donc une base de \(M_{n,p}(\K)\), contenant \(np\) éléments. Ainsi
\[
\dim(M_{n,p}(\K))=np.
\]
\end{proof}

\begin{remarque}
Cette base \((E_{ij})\) est la base canonique de l’espace des matrices.
\end{remarque}

\section{Méthodes pratiques}

\begin{methode}
Pour additionner ou multiplier une matrice par un scalaire, on travaille coefficient par coefficient.
\end{methode}

\begin{methode}
Pour calculer un produit matriciel :
\begin{enumerate}
    \item vérifier d’abord que le produit a un sens ;
    \item déterminer la taille de la matrice résultat ;
    \item calculer chaque coefficient comme produit d’une ligne par une colonne.
\end{enumerate}
\end{methode}

\begin{methode}
Pour montrer qu’une matrice \(B\) est l’inverse de \(A\), il faut vérifier
\[
AB=BA=I_n.
\]
\end{methode}

\begin{methode}
Pour résoudre un système carré \(AX=B\), si \(A\) est inversible, penser immédiatement à
\[
X=A^{-1}B.
\]
\end{methode}

\section{Erreurs classiques}

\begin{itemize}
    \item Oublier de vérifier que le produit matriciel a un sens.
    \item Croire que deux matrices de tailles différentes peuvent s’additionner.
    \item Croire que le produit matriciel est commutatif.
    \item Oublier que \((AB)^T=B^TA^T\).
    \item Écrire à tort
    \[
    (AB)^{-1}=A^{-1}B^{-1}.
    \]
    La bonne formule est
    \[
    (AB)^{-1}=B^{-1}A^{-1}.
    \]
    \item Croire que \(AB=0\) implique \(A=0\) ou \(B=0\).
    \item Utiliser la formule de l’inverse d’une matrice \(2\times 2\) sans vérifier que \(ad-bc\neq 0\).
\end{itemize}

\section{Résumé du chapitre}

\begin{itemize}
    \item Une matrice est un tableau de coefficients.
    \item L’ensemble \(M_{n,p}(\K)\) est un espace vectoriel de dimension \(np\).
    \item Le produit matriciel est défini lorsque les tailles sont compatibles.
    \item Il est associatif et distributif, mais en général non commutatif.
    \item La transposée échange les lignes et les colonnes, et vérifie
    \[
    (AB)^T=B^TA^T.
    \]
    \item Une matrice carrée \(A\) est inversible s’il existe \(A^{-1}\) tel que
    \[
    AA^{-1}=A^{-1}A=I_n.
    \]
    \item Si \(A\) est inversible, alors le système
    \[
    AX=B
    \]
    admet une unique solution
    \[
    X=A^{-1}B.
    \]
    \item Les opérations élémentaires sur les lignes correspondent à des multiplications par des matrices élémentaires.
\end{itemize}

\section*{Exercices d’application immédiate}

\begin{enumerate}[label=\textbf{Exercice \arabic*.}]
    \item Déterminer la taille des matrices suivantes :
    \[
    \begin{pmatrix}
    1 & 2 & 3\\
    4 & 5 & 6
    \end{pmatrix},
    \qquad
    \begin{pmatrix}
    1\\
    2\\
    3
    \end{pmatrix},
    \qquad
    \begin{pmatrix}
    1 & 0\\
    0 & 1
    \end{pmatrix}.
    \]

    \item Calculer :
    \[
    \begin{pmatrix}
    1 & 2\\
    3 & 4
    \end{pmatrix}
    +
    \begin{pmatrix}
    0 & -1\\
    5 & 2
    \end{pmatrix}.
    \]

    \item Calculer :
    \[
    3
    \begin{pmatrix}
    1 & -2\\
    0 & 4
    \end{pmatrix}.
    \]

    \item Calculer les produits suivants lorsqu’ils ont un sens :
    \[
    \begin{pmatrix}
    1 & 2\\
    3 & 4
    \end{pmatrix}
    \begin{pmatrix}
    5 & 6\\
    7 & 8
    \end{pmatrix},
    \qquad
    \begin{pmatrix}
    1 & 2 & 3
    \end{pmatrix}
    \begin{pmatrix}
    1\\
    0\\
    -1
    \end{pmatrix}.
    \]

    \item Calculer la transposée de
    \[
    \begin{pmatrix}
    1 & 2 & 3\\
    4 & 5 & 6
    \end{pmatrix}.
    \]

    \item Vérifier sur un exemple que, en général,
    \[
    AB\neq BA.
    \]

    \item Déterminer l’inverse de
    \[
    \begin{pmatrix}
    1 & 2\\
    3 & 4
    \end{pmatrix}.
    \]

    \item Écrire sous forme matricielle le système
    \[
    \begin{cases}
    x+2y-z=3,\\
    2x-y+4z=1.
    \end{cases}
    \]

    \item Résoudre le système
    \[
    \begin{cases}
    x+y=1,\\
    2x+3y=4
    \end{cases}
    \]
    par la méthode de l’inverse de matrice.

    \item Montrer que les matrices élémentaires traduisent les opérations élémentaires sur les lignes.
\end{enumerate}

\end{document}